\section{The combinatorial scan, and what has never been looked at}
\label{sec:combinatorial}

Section~\ref{sec:budget} priced the program that exists: the spectra published searches scan. A
scan does not have to be restricted to those. If an analysis builds every invariant mass it can
from the objects in an event, the relevant budget is a different and much larger one, and the
difference between the two is precisely the space nobody has examined.

\subsection{Budget of a fully combinatorial scan}

To keep this concrete rather than hypothetical we fix one object budget and price it. It is the
shape a data-directed scan~\cite{ddp} of this kind naturally takes, and it is a design choice of
this note rather than a description of any particular implementation:

\begin{itemize}[leftmargin=1.4em,itemsep=2pt]
  \item Categories are \textbf{exclusive} multiplicity bins in
        $(n_e, n_\mu, n_j, n_b, n_Z, \mathrm{MET})$, so different categories select disjoint event
        sets and are statistically independent looks.
  \item The object budget is $n_e \le 2$, $n_\mu \le 2$, $n_j \le 3$, $n_b \le 3$, $n_Z \le 1$,
        with $2 \le \sum n \le 4$, missing energy excluded from the masses.
  \item A single-lepton trigger requires each populated category to contain a lepton, either a
        loose $e/\mu$ or a leptonic $Z$.
  \item Missing energy splits every category in two but never enters a mass, and two same-flavour
        loose leptons split into opposite- and same-sign.
  \item Every size-2 to size-4 subset of the indexed objects is its own histogram, so $j_0j_1$,
        $j_0j_2$ and $j_1j_2$ are three spectra rather than one.
\end{itemize}

This yields \textbf{150 exclusive multiplicity categories}, which become 204 once the
same-flavour dilepton categories are split into opposite- and same-sign, and \textbf{1094
(category, mass-group) combinations}, or 1502 histograms after that same charge split. They span 82
distinct flavour compositions. Assigning each spectrum a resolution from its object
composition ($\sigma/p_{\mathrm{T}} = 0.04$ for $e/\mu/Z$, $0.10$ for light jets and $0.20$ for
$b$-jets, combined as $r = \tfrac12\sqrt{\langle \sigma_i^2\rangle}$, which reproduces the
per-channel values used in Section~\ref{sec:budget}) and applying Eq.~\eqref{eq:ns} gives
\begin{equation*}
  \Nsig = 1.43\times10^{5}
  \qquad\Longrightarrow\qquad
  \Zloc = 6.98 \quad (\text{band } 6.88\text{--}7.08),
\end{equation*}
so a local $5\sigma$ in such a scan is worth only $1.1\sigma$ globally. Relative to the
published-program budget the trials factor grows by a factor 39 while the discovery bar moves by
$0.5\sigma$, which is the logarithm doing its work again. A combinatorial scan is not statistically
prohibitive. It is about half a sigma more expensive than the program ATLAS already runs.

\subsection{Which compositions has ATLAS actually scanned?}

Those same 82 compositions give a natural denominator for a coverage question that is normally
asked qualitatively. We map them onto the published bump observables of Appendix~\ref{app:windows}
generously, counting a composition as covered if \emph{any} published ATLAS search scans a mass
built from those object types. A published ``$V$'' is allowed to be a $Z$, a $jj$ pair or an
$\ell\ell$ pair, a published ``$t$'' to contain $b+$jets, and so on. Even on that permissive
reading:

\begin{table}[htbp]
\centering
\caption{Reachable ($\le4$-object) mass compositions with and without a published ATLAS bump hunt.
The mapping is deliberately generous, so the uncovered column is a lower bound on the gap.}
\label{tab:compgap}
\begin{tabular}{lrrr}
\toprule
multiplicity & compositions & with a published scan & with none \\
\midrule
$K=2$ & 14 & 14 & 0 \\
$K=3$ & 28 & 7 & 21 \\
$K=4$ & 40 & 4 & 36 \\
\midrule
total & 82 & 25 & \textbf{57} \\
\bottomrule
\end{tabular}
\end{table}

In this basis the gap is a pure multiplicity effect. Two-body space is saturated, while 75\,\% of
three-body and 90\,\% of four-body compositions have never been scanned. An independent audit built
from the publication record upward, in a different object basis
$\{$light jet, $b$-jet, lepton, photon$\}$, reproduces the shape: 90\,\% of two-body, 25\,\% of
three-body and roughly 11\,\% of four-body combinations scanned.

\subsection{Three reasons the covered fraction is overstated}

The generous mapping hides structure that sharpens the conclusion considerably.

\paragraph{Tag versus scan.} Several searches that look like higher-multiplicity coverage in fact
scan a two-body mass, with the extra object used only as a trigger or flavour tag. The $b$-tagged
dijet search~\cite{atlas_btag_dijet} is explicitly a two-body scan, and the dijet$+$lepton
search~\cite{atlas_dijet_lepton} scans $m_{jj}$ only. Counting either as multi-body coverage
overstates the breadth of the record.

\paragraph{Almost all mixed-flavour coverage is three papers, all since 2022.} Of roughly 45
distinct spectra ever bump-hunted by ATLAS, only about 15 are mixed-flavour or have three or more
bodies, and essentially all of those come from~\cite{atlas_ml_anomaly, atlas_multibody,
atlas_boostedz}. One of the three is a machine-learning anomaly search; the other two are
conventional scans that were simply never run before. Before 2022 that coverage was close to zero.

\paragraph{Most covered two-body spectra are covered only inside anomaly-selected phase space.}
The broadest mixed-flavour scan in the ATLAS record~\cite{atlas_ml_anomaly} covers $m_{jb}$,
$m_{je}$, $m_{be}$, $m_{j\mu}$, $m_{b\mu}$, $m_{j\gamma}$ and $m_{b\gamma}$ from $0.3$ to $8\TeV$,
but does so inside autoencoder-selected regions rather than inclusive selections. There is no
inclusive dedicated search for $m(\ell b)$, $m(\gamma b)$, or non-leptoquark $m(\ell j)$.

\subsection{The programme that was meant to cover the rest}

The ATLAS general search~\cite{atlas_general_search} populates 704 event classes, analyses 686 of
them and defines more than $10^5$ regions. It reads like blanket coverage. It is not, for a reason
that bears directly on Section~\ref{sec:budget}: it scans exactly two event-level observables per
class, $m_{\mathrm{eff}}$ and the invariant mass of \emph{all} visible objects. Class ``$1e\,2j$''
therefore yields $m_{\mathrm{inv}}(e,j,j)$ and $m_{\mathrm{eff}}$, never $m(ej)$ or $m(jj)$. The
general search is structurally blind to sub-combination masses, which is exactly where resonances
live.

The paper states the reason outright: the observable set was restricted ``to avoid a large
increase in the number of hypothesis tests''. \textbf{ATLAS declined this coverage explicitly, on
look-elsewhere grounds.} That is a deliberate trade rather than an oversight, and it is precisely
the trade the numbers above reprice. Moving from two event-level observables to all
sub-combinations shifts the discovery bar from $6.4\sigma$ to $7.0\sigma$. That is the entire cost
of the coverage the general search gave up.

Two further points on that programme. It ran on $3.2\fbinv$, which is 2015 data and 2.3\,\% of
Run~2, and it has had no successor in the event-class format. The full-Run-2 anomaly-detection
scan~\cite{atlas_ml_anomaly} is a successor in spirit, but it covers nine two-body spectra rather
than hundreds of exclusive classes.

\subsection{Named gaps}

Concrete negatives, each checked individually against the publication record rather than inferred
from the table above:

\begin{itemize}[leftmargin=1.4em,itemsep=2pt]
  \item \textbf{$m(\ell\gamma)$}, the only uncovered two-body composition in the basis. Excited
        leptons $\ell^\ast \to \ell\gamma$ were searched for at $7\TeV$~\cite{atlas_lstar_7tev} and
        $8\TeV$~\cite{atlas_lstar_8tev}; every $13\TeV$ excited-lepton search uses other channels.
  \item \textbf{Same-sign dilepton plus jets mass spectra.} $H^{\pm\pm} \to \ell^\pm\ell^\pm$ was a
        genuine same-sign mass scan at $7\TeV$~\cite{atlas_dch_7tev}, but the full-Run-2
        update~\cite{atlas_dch_run2} replaced the scan with a counting signal region.
  \item \textbf{$m(jj\gamma)$, $m(bbj)$, $m(\ell\ell j)$}, clean three-body negatives.
  \item \textbf{$m(jjjj)$}: no light-jet four-body scan. The $b$-tagged analogue \emph{is} covered,
        by the resolved $X\to HH \to 4b$ channel~\cite{atlas_hh4b}.
  \item \textbf{Colour-sextet and $E_6$ diquarks}, \textbf{$b^\ast \to b\gamma$} (CMS sets excited-bottom-quark limits in its $\gamma+$jet resonance
search~\cite{cms_gammajet}; the ATLAS $\gamma+$jet searches are flavour-inclusive), \textbf{$\tau^\ast \to \tau\gamma$} (the $\tau^\ast$ \emph{model} was
        searched through contact interactions~\cite{atlas_taustar}, so the $\tau\gamma$
        \emph{channel} is the gap), \textbf{$m(t\gamma)$}, and \textbf{same-sign $t\bar t$ as a mass
        bump}~\cite{atlas_ss_tt}: model or channel gaps with no ATLAS scan in any run.
  \item \textbf{The entire mixed-flavour $K=3$ and $K=4$ block}: $m_{jb\mu}$, $m_{ejj}$,
        $m_{e\mu b}$, $m_{eejb}$, $m_{\mu\mu bb}$ and the rest.
\end{itemize}

A fourth category turned up that is arguably more interesting than any of these, and is not
``never examined'': signatures examined once at Run-1 luminosity and then quietly abandoned.
$m(\ell Z)$~\cite{atlas_vll_8tev} and $m(Zb)$~\cite{atlas_vlq_zb_8tev} were both scanned at
$8\TeV$ and never repeated. Re-running an abandoned $8\TeV$ scan on the full Run-2 dataset costs
nothing in trials that Section~\ref{sec:budget} has not already counted.

A closing caution on how to read all of this. ``Never scanned'' is a statement about the
publication record, not about whether the signal models exist or whether samples were ever
generated for them. Those are orthogonal axes: a channel can be well supplied with theory and
simulation and still have no published scan, which is precisely the situation the table above
counts.
